% !TeX root = ../教材题库_全部条目_转写工作稿.tex
% Source: XJ-B1-C05-HW-P202-Y0993
% 第5章 / 5.5 三角函数模型的简单应用 / 教材 p.202 / PDF 第209页
% OCR未校对：公式与文字均可能有误；未完成数学TeX转写；禁止直接发给学生
% 原题及相邻原解截图（核对用）：../题目原图/XJ-B1-C05-HW-P202-Y0993-01.webp
% 本文件只在转写完成后作为题干使用；OCR 条目尚不是可靠题干。
\begin{bankitem}{习题5.5 第5题}{第5章 / 5.5 三角函数模型的简单应用 / 教材 p.202 / PDF 第209页}{XJ-B1-C05-HW-P202-Y0993}{OCR未校对：公式与文字均可能有误；未完成数学TeX转写；禁止直接发给学生}
\begin{ocrdraft}
[教材p.202 / PDF第209页]
\par
5．已知某海滨浴场的浪高3，（m）是时间底时）（0<
\par
0
\par
8
\par
14
\par
对时
\par
／24）的函数，记作y=f(t).下表是某日各时刻的浪
\par
（第4题）
\par
高数据：
\par
0时
\par
0
\par
3
\par
6
\par
9
\par
12
\par
15
\par
18
\par
21
\par
24
\par
y/m
\par
0．99
\par
经长期观测，一70）可近似地看成是函数y=Acos十b\ensuremath{\cdot}
\par
（1）根据以上数据，求出该函数的周期T、振幅A及函数解析式；
\par
（2）依据规定，当海浪高度高于1m时才对冲浪爱好者开放，试依据（1）的结
\par
论，判断一天内8：00至20：00之间有多长时间可供冲浪者进行运动．
\end{ocrdraft}
\end{bankitem}
